\documentclass[11pt]{ctexart}
\usepackage[a4paper,margin=1in]{geometry}
\usepackage{graphicx}
\usepackage{xcolor}
\usepackage{hyperref}
\definecolor{refgray}{gray}{0.45}
\title{\textbf{市场最新观点汇总}\\[0.4em]{\large 全球资金正从分散试探转向集中押注科技与AI，利率市场定价逻辑从通胀冲击转向央行反应函数，新兴市场出现结构性分化。}}
\author{KC桌面 自动生成}
\date{260620}
\begin{document}
\maketitle
\tableofcontents
\newpage
\section{一页摘要}
\begin{itemize}
  \item 全球利率市场正经历从通胀冲击定价向央行反应函数定价的转变，利率波动率重心从长端向短端转移，曲线中段（5年期）成为风险收益比最优的区间。
  \item 全球资金流向显示，科技板块吸金能力加速，AI驱动的需求不仅推高增长预期，还通过推高通胀和中性利率预期强化美元强势，科技板块已成为全球宏观资产定价的核心锚点。
  \item 韩国市场短期动能来自事件驱动和外资回补，但主动基金仓位仍在下降，中期结构性折价的修复取决于MSCI评估中制度性障碍能否真正落地。
  \item 全球资本开支周期正在从规模扩张转向议价权重构，AI数据中心、油气勘探和电力基础设施是三大受益赛道，但个股分化将显著加大。
  \item 期限溢价模型更新揭示，传统收益率曲线模型可能系统性高估当前实际期限溢价水平，调查增强模型显示美国10年期期限溢价约70bp，比传统模型低40bp。
\end{itemize}
\section{宏观与利率：央行反应函数重塑曲线定价}
\textbf{市场对通胀的敏感度正从能源冲击转向央行沟通机制，利率波动率重心从长端向短端转移，曲线中段提供不对称机会。}

\begin{itemize}
  \item GS指出，美伊协议带来的能源缓解已部分消化，取而代之的是美联储鹰派沟通对短端利率的重塑，市场正从宏观冲击主导转向政策不确定性主导的环境。
  \item OIS曲线显示市场几乎定价了两次加息，但利率曲线的峰值并未显著高于5月非农数据后的水平，形成鹰派前端与丰富长端的格局。
  \item GS认为5年期国债在当前环境下风险收益比最优：若加息风险落地，5年相对2年和10年会跑赢；若数据继续走强，5年同样受益。
  \item 传统收益率曲线模型可能系统性高估当前实际期限溢价水平，因为它未能区分结构性利率预期上移和周期性风险定价。
  \item 调查增强模型估算美国10年期期限溢价约70bp，比传统模型低40bp；欧洲低约15bp；英国是G4中最低的，说明英国利率预期本身很高。
  \item 日本是两种模型差异最极端的案例：传统模型将收益率曲线陡峭化几乎全归因于期限溢价上升，但调查模型显示很大一部分来自政策利率路径预期的结构性抬升。
  \item G4期限溢价上升与央行缩表同步，但近期隐含波动率下降尚未反映在期限溢价中，意味着期限溢价仍有回落空间。
  \item 欧洲方面，油价缓解降低了利率波动，利好主权债券利差收窄，但能源期货已低于研报预期，继续利好空间有限。
  \item 英国通胀数据走弱叠加失业率趋势下行，央行大概率按兵不动，市场定价的45bp加息预期有下行空间。
\end{itemize}
\begin{figure}[htbp]
\centering
\includegraphics[width=0.88\linewidth]{figures/fig_001.jpg}
\caption{Exhibit 2 - Exhibit 1: Hike risk is more front-loaded, but the curve is not pricing a significantly higher peak 3m SOFR futures curve}
\end{figure}
\begin{figure}[htbp]
\centering
\includegraphics[width=0.88\linewidth]{figures/fig_002.jpg}
\caption{Exhibit 2 - Exhibit 2: 5s should outperform on the curve from front-loading of hike risk against stickier or declining forwards further out the curve Shift in curve segments implied by changes in pace of hikes over next 6m, extent of hikes ove}
\end{figure}
\begin{figure}[htbp]
\centering
\includegraphics[width=0.88\linewidth]{figures/fig_004.jpg}
\caption{Exhibit 4 - Exhibit 4: Greater policy uncertainty flattens the volatility tail curve 1y3m US rate uncertainty (orthogonalized against inflation, growth) vs implied vol tail curve}
\end{figure}
\begin{figure}[htbp]
\centering
\includegraphics[width=0.88\linewidth]{figures/fig_005.jpg}
\caption{Exhibit 5 - Exhibit 5: Markets are hawkish vs policy rate surveys ECB surveys of the policy rate vs market probabilities based on swaption pricing}
\end{figure}
\begin{figure}[htbp]
\centering
\includegraphics[width=0.88\linewidth]{figures/fig_007.jpg}
\caption{Exhibit 1 - Exhibit 1: Survey-based estimations currently sit below yield curve only term premium estimates ...}
\end{figure}
\begin{figure}[htbp]
\centering
\includegraphics[width=0.88\linewidth]{figures/fig_008.jpg}
\caption{Exhibit 2 - Exhibit 2: ... and also show relatively more stability}
\end{figure}
\begin{figure}[htbp]
\centering
\includegraphics[width=0.88\linewidth]{figures/fig_009.jpg}
\caption{Exhibit 1 - Exhibit 3: UK survey-based estimates are the lowest among the G4, indicating an important part comes from high rates expectations}
\end{figure}
\begin{figure}[htbp]
\centering
\includegraphics[width=0.88\linewidth]{figures/fig_010.jpg}
\caption{Exhibit 4 - Exhibit 4: Using surveys shows that the rise in JGBs is not only a term premium story}
\end{figure}
\begin{figure}[htbp]
\centering
\includegraphics[width=0.88\linewidth]{figures/fig_012.jpg}
\caption{Exhibit 6 - Exhibit 6: Our existing estimates of term premium are more volatile than survey-augmented estimates Standard Deviation of Term Premium Estimates}
\end{figure}
\begin{figure}[htbp]
\centering
\includegraphics[width=0.88\linewidth]{figures/fig_018.jpg}
\caption{Exhibit 12 - Exhibit 12: G4 term premium has risen alongside shrinking central bank balance sheets Average of yield curve-only and survey-based term premium across G4 (GDP weighted) and average central bank holdings (GDP weighted)}
\end{figure}
\begin{figure}[htbp]
\centering
\includegraphics[width=0.88\linewidth]{figures/fig_019.jpg}
\caption{Exhibit 13 - Exhibit 13: The recent decline in implied volatility has not yet been reflected in lower term premium Average of yield curve-only and survey-based term premium across G4 (GDP weighted) and average implied vol (GDP weighted)}
\end{figure}
{\small\color{refgray}\textbf{References:} [R001] GS：市场低估的不是通胀本身，而是通胀预期的再定价机制; [R002] GS：市场真正低估的不是期限溢价的水平，而是它为何不再回落}

\section{AI与算力：资本开支重塑行业竞争格局}
\textbf{全球进入资本开支驱动的后现代周期，回报逻辑从规模扩张转向议价权重构，AI基础设施和电力公用事业是核心受益者。}

\begin{itemize}
  \item GS全球策略师Peter Oppenheimer指出，当前资本开支周期回报更多来自每股盈利增长而非估值扩张，个股分化将显著加大，主动选股空间正在打开。
  \item 欧洲数据中心管道已接近500GW，相当于当前欧盟电力需求的1.5倍，六个月前仅约300GW，增长70\%。
  \item 超大规模云服务商的AI基础设施需求依然强劲，数据中心定价权持续增强，半导体设备和电力电网相关公司成为直接受益者。
  \item 自主软件开发正在催生暗工厂模式：编码AI写代码，运维AI自动诊断和修复问题，形成更紧密的开发-运维闭环。
  \item 尽管油价回落至80美元/桶附近，石油公司的资本开支正在回暖，原因是过去10年储备寿命下降了60\%，美国页岩油对油价的压制效应正在消退。
  \item 能源安全压力进一步推动前沿勘探回归，某地震勘探公司反馈EAGE行业会议参会人数创纪录。
  \item 资本开支的准入门槛正在提高，只有拥有稳定现金流、明确战略方向和技术壁垒的企业才能持续投入并看到回报。
  \item 谁能在未来几年获得稳定的电力供应、在电网接入和土地审批上占据先机，谁就能在数据中心竞争中建立壁垒。
\end{itemize}
{\small\color{refgray}\textbf{References:} [R005] GS：市场低估的不是资本开支本身，而是资本开支正在重塑行业竞争格局}

\section{能源与大宗：资本开支回暖与结构性变化}
\textbf{石油资本开支在储备寿命下降和能源安全压力下回归，但油价回落限制了上行空间；欧洲能源缓解利好利差收窄，但空间有限。}

\begin{itemize}
  \item 石油公司资本开支回暖，尽管油价回落至80美元/桶附近，但过去10年储备寿命下降60\%是核心驱动。
  \item 美国页岩油对油价的压制效应正在消退，能源安全压力进一步推动前沿勘探回归。
  \item 欧洲方面，油价缓解降低了利率波动，利好主权债券利差收窄，但能源期货已低于研报预期，继续利好空间有限。
  \item 欧洲数据中心管道快速增长，电力基础设施将迎来长期需求支撑，公用事业板块受益。
  \item 资本开支的回报逻辑从规模扩张转向议价权重构，能够将资本投入转化为稀缺资源控制权的企业是赢家。
\end{itemize}
{\small\color{refgray}\textbf{References:} [R001] GS：市场低估的不是通胀本身，而是通胀预期的再定价机制; [R005] GS：市场低估的不是资本开支本身，而是资本开支正在重塑行业竞争格局}

\section{区域市场：韩国结构性分歧与MSCI评估}
\textbf{韩国市场短期动能来自事件驱动和外资回补，但主动基金仓位持续下降，中期折价修复取决于MSCI制度性改革落地。}

\begin{itemize}
  \item KOSPI一周暴涨11\%突破9000点，科技股领涨，外资净买入2.5万亿韩元，但上涨广度有限。
  \item EPFR数据显示，亚洲基金和新兴市场基金对韩国股票仍显著低配，亚洲基金低配195bp，新兴市场基金低配285bp。
  \item 过去一个月亚洲基金减持韩国105bp，过去三个月累计减持140bp；新兴市场基金分别减持165bp和195bp。
  \item 本周外资净买入更可能来自交易型或被动型资金短期回补，而非主动管理型基金的战略性加仓。
  \item MSCI评估中韩国仍有5项需改善评级：外汇自由化、投资者注册、英文信息披露、清算结算、可转移性。
  \item MSCI表示要看到实际改善而非仅凭改革承诺，6月23日年度评估大概率不会纳入发达市场观察名单。
  \item 保险(+18.4\%)、科技(+16.6\%)、零售(+5.4\%)领跑，建筑(-9.9\%)、电信(-6.2\%)、休闲(-5.1\%)垫底。
  \item 12个月前瞻EPS上调1.5\%，科技板块盈利预期最强，化工板块下调最多。
  \item 韩元兑美元贬值0.8\%，韩国权益风险指标回到中性区域。
\end{itemize}
\begin{figure}[htbp]
\centering
\includegraphics[width=0.88\linewidth]{figures/fig_020.jpg}
\caption{Exhibit 3 - Exhibit 3: Among EM \& Asian focused funds, Insurance and Tech H/W are the most OW while Internet and Health Care are the most UW Korea allocation in EM \& Asian focused funds (OW/UW vs. benchmark, bp)}
\end{figure}
\begin{figure}[htbp]
\centering
\includegraphics[width=0.88\linewidth]{figures/fig_021.jpg}
\caption{Exhibit 4 - Exhibit 4: Over the past 3 months, EM \& Asian funds have raised their exposure to Industrials the most while trimmed down Tech H/W exposure by sector Change in Korea allocation in EM \& Asian focused funds (past 3M, bp)}
\end{figure}
\begin{figure}[htbp]
\centering
\includegraphics[width=0.88\linewidth]{figures/fig_030.jpg}
\caption{Exhibit 16 - Exhibit 16: Performance of Korean and regional equity markets Exhibit 17: KOSPI and MSCI regional index KOSPI Index price performance (KRW)}
\end{figure}
\begin{figure}[htbp]
\centering
\includegraphics[width=0.88\linewidth]{figures/fig_031.jpg}
\caption{Exhibit 18 - Exhibit 18: Weekly sector performance relative to KOSPI}
\end{figure}
\begin{figure}[htbp]
\centering
\includegraphics[width=0.88\linewidth]{figures/fig_032.jpg}
\caption{Exhibit 19 - Exhibit 19: Historical trend of monthly KOSPI earnings momentum}
\end{figure}
\begin{figure}[htbp]
\centering
\includegraphics[width=0.88\linewidth]{figures/fig_033.jpg}
\caption{Exhibit 20 - Exhibit 20: Weekly earnings momentum for KOSPI and selected sectors}
\end{figure}
\begin{figure}[htbp]
\centering
\includegraphics[width=0.88\linewidth]{figures/fig_034.jpg}
\caption{Exhibit 22 - Exhibit 23: 12-month forward P/E and P/B for KOSPI, since Jan 2008}
\end{figure}
\begin{figure}[htbp]
\centering
\includegraphics[width=0.88\linewidth]{figures/fig_037.jpg}
\caption{Exhibit 27 - Exhibit 27: MXKR valuation discount relative to MSCI AC World NTM P/E}
\end{figure}
\begin{figure}[htbp]
\centering
\includegraphics[width=0.88\linewidth]{figures/fig_038.jpg}
\caption{Exhibit 28 - Exhibit 28: MXKR valuation discount relative to MXAPJ NTM P/E}
\end{figure}
\begin{figure}[htbp]
\centering
\includegraphics[width=0.88\linewidth]{figures/fig_044.jpg}
\caption{Exhibit 39 - Exhibit 39: USDKRW and USDCNY rates}
\end{figure}
\begin{figure}[htbp]
\centering
\includegraphics[width=0.88\linewidth]{figures/fig_045.jpg}
\caption{Exhibit 40 - Exhibit 40: USDKRW vs. 10 year bond yields differentials between the US and Korea}
\end{figure}
\begin{figure}[htbp]
\centering
\includegraphics[width=0.88\linewidth]{figures/fig_046.jpg}
\caption{Exhibit 41 - Exhibit 41: Korea Equity Risk Barometer}
\end{figure}
{\small\color{refgray}\textbf{References:} [R003] GS：韩国市场上涨背后的结构性分歧，比指数本身更值得关注}

\section{全球资金流向：科技虹吸与新兴市场分化}
\textbf{全球资金正从分散试探转向集中押注科技与AI，中国内地股票基金持续失血，新兴市场出现明显分化。}

\begin{itemize}
  \item 截至6月17日当周，全球股票基金净流入1260亿美元，远超此前一周的310亿美元。
  \item 科技板块和工业板块在全球范围内录得最大净流入，美国科技基金流入尤为强劲，AI驱动的需求是背后关键因素。
  \item GS指出，强劲的AI需求不仅推高增长预期，还通过推高通胀和中性利率预期强化美元强势，科技板块已成为全球宏观资产定价的核心锚点。
  \item 全球固收基金净流入191.6亿美元，短久期和通胀保护债券持续受欢迎。
  \item 新兴市场出现分化：中国大陆权益基金净流出，台湾净流入；硬通货债券基金净流入，本币债券基金净流出。
  \item 跨境外汇流中，美元、欧元、日元需求最强，人民币净流出最大。
  \item 年初至今累计数据看，工业、基础设施、能源板块资金流入最多，科技板块累计流入并不算最突出，消费和金融板块净流出。
  \item 台湾、巴西、韩国是资金流入最多的新兴市场，中国大陆和印度净流出最多。
\end{itemize}
{\small\color{refgray}\textbf{References:} [R004] GS：资金流向正在重塑全球资产定价的底层逻辑}

\section{风险偏好与市场情绪：权益风险指标回归中性}
\textbf{韩国市场情绪指标显示风险偏好回归中性，但主动基金仓位与指数上涨的背离值得警惕。}

\begin{itemize}
  \item 韩国权益风险指标回到中性区域，VKOSPI波动率指数处于低位。
  \item KOSPI中高于200日移动均线的成分股比例回升，但上涨广度有限。
  \item 外资净买入主要流向科技和汽车板块，但主动基金仓位仍在下降，形成方向性背离。
  \item 韩国保险和科技硬件最受超配，互联网和医疗最受低配。
  \item 货币市场基金资产增加250亿美元，显示部分资金仍保持谨慎。
  \item 全球资金从部分新兴市场撤出转向美元资产，美元、欧元、日元需求最强。
\end{itemize}
\begin{figure}[htbp]
\centering
\includegraphics[width=0.88\linewidth]{figures/fig_046.jpg}
\caption{Exhibit 41 - Exhibit 41: Korea Equity Risk Barometer}
\end{figure}
\begin{figure}[htbp]
\centering
\includegraphics[width=0.88\linewidth]{figures/fig_048.jpg}
\caption{Exhibit 43 - Exhibit 43: Korea Volatility Index (VKOSPI)}
\end{figure}
\begin{figure}[htbp]
\centering
\includegraphics[width=0.88\linewidth]{figures/fig_049.jpg}
\caption{Exhibit 44 - Exhibit 44: KOSPI Percent of Members above 200 day moving average}
\end{figure}
{\small\color{refgray}\textbf{References:} [R003] GS：韩国市场上涨背后的结构性分歧，比指数本身更值得关注; [R004] GS：资金流向正在重塑全球资产定价的底层逻辑}

\section{结语}
综合来看，当前市场主线围绕央行反应函数重塑、AI资本开支周期和全球资金再配置三条线索展开，利率曲线中段、科技基础设施和具备制度性改革潜力的新兴市场是值得持续关注的领域。
\newpage
\section{报告覆盖清单}
\begin{itemize}
  \item [R001] 宏观与利率：央行反应函数重塑曲线定价 / 能源与大宗：资本开支回暖与结构性变化 - GS：市场低估的不是通胀本身，而是通胀预期的再定价机制
  \item [R002] 宏观与利率：央行反应函数重塑曲线定价 - GS：市场真正低估的不是期限溢价的水平，而是它为何不再回落
  \item [R003] 区域市场：韩国结构性分歧与MSCI评估 / 风险偏好与市场情绪：权益风险指标回归中性 - GS：韩国市场上涨背后的结构性分歧，比指数本身更值得关注
  \item [R004] 全球资金流向：科技虹吸与新兴市场分化 / 风险偏好与市场情绪：权益风险指标回归中性 - GS：资金流向正在重塑全球资产定价的底层逻辑
  \item [R005] AI与算力：资本开支重塑行业竞争格局 / 能源与大宗：资本开支回暖与结构性变化 - GS：市场低估的不是资本开支本身，而是资本开支正在重塑行业竞争格局
\end{itemize}
\section{逐篇报告摘录}
\subsection{[R001] GS：市场低估的不是通胀本身，而是通胀预期的再定价机制}
{\small\color{refgray}归类：宏观与利率：央行反应函数重塑曲线定价 / 能源与大宗：资本开支回暖与结构性变化}

这份GS全球利率策略报告的核心判断，并非关于通胀是否见顶，而是关于一个更微妙、也更具交易含义的转变：通胀风险已经从单纯的能源价格冲击，转移到了央行反应函数的重新定价上。这意味着，市场对通胀的敏感度正在发生结构性变化，投资者需要调整的不是方向判断，而是对利率曲线不同段落的定价逻辑。

报告明确指出，美伊协议带来的能源缓解已经部分消化，但取而代之的是美联储鹰派沟通对短端利率的重塑。GS团队认为，市场正在从一个以宏观冲击为主导的环境，转向一个以政策不确定性为主导的环境。这个转变的直接后果是：利率波动率的重心将从长端向短端转移，而曲线中段（即所谓的“belly”）可能成为风险收益比最优的区间。

这不是一个简单的“通胀见顶，做多债券”的故事。这是一个关于波动率结构、央行沟通机制和曲线形态重新定价的深度叙事。以下，我们从五个层次展开这份报告的核心洞察。

1. 鹰派预期的前置化正在压缩长端风险溢价，但曲线中段反而提供了不对称机会

GS团队观察到，市场对加息的定价已经从前期的“分散化”转向了“前置化”。OIS曲线显示，市场几乎定价了两次加息，但有趣的是，利率曲线的峰值并未显著高于5月非农数据出炉后的水平。这意味着什么？

这意味着市场对美联储反应函数的理解正在发生变化。过去，通胀数据超预期会同时推升短端和长端利率；但现在，市场似乎认为，只要美联储展现出足够的前瞻性，长端的风险溢价就会被压制。GS将这种现象描述为“鹰派的前端/丰富的长端”格局。这种格局在历史上往往对应着曲线陡峭化的倾向。

但GS团队并未简单地推荐陡峭化交易。他们认为，在短期内，市场对反应函数的重新评估更有利于曲线中段（5年期）的表现。逻辑在于：如果加息风险被前置化，那么5年期相对于2年期和10年期，在曲线上的相对价值会提升。报告中的Exhibit 2清晰地展示了这一机制：当加息节奏加快但终端利率不变时，2s5s段会趋陡，而5s10s段会趋平。这意味着，5年期在一个典型的“鹰派前置”场景中，是曲线上的“甜点”。

*So what**: 对于利率投资者而言，单纯做陡或做平曲线可能都不是最优解。GS推荐的“做多曲线中段”策略，本质上是在押注一个特定的波动率结构：短端的波动率上升，但长端的风险溢价被压制。这是一个需要精细调整久期和凸性的交易，而非方向性押注。

这是报告中最具洞察力的部分之一。GS团队指出，能源价格下跌和地缘政治不确定性降低，已经压低了整个波动率曲面的水平。但与此同时，美联储反应函数的不确定性正在上升。这种从“宏观波动”到“政策波动”的切换，对波动率曲线的形态有明确的含义。

GS的波动率模型显示，当通胀调查的离散度下降、但政策利率调查的离散度上升时，波动率尾部曲线（即长期波动率与短期波动率之差）通常会趋平。Exhibit 4的散点图清晰地展示了这一负相关关系。这意味着，短期利率的波动率可能会上升，但长期利率的波动率反而会相对稳定。

这个结论对于利率期权交易者至关重要。它意味着，押注长期波动率上升的策略（如做多长期跨式期权）可能面临不利环境，而做空尾部风险或做多短期波动率的策略可能更符合当前的市场结构。GS团队特别提到，随着市场适应美联储的新沟通机制，FOMC会议和会议纪要发布前后的局部波动率仍将活跃。

*So what**: 波动率交易的核心逻辑正在发生变化。过去两年，投资者习惯于用长期波动率来对冲地缘政治和通胀风险。但未来，波动率的主要来源将是央行的沟通和反应函数。这意味着，波动率交易的周期将缩短，事件驱动型策略将比趋势型策略更有效。

欧洲部分的分析是报告的另一大亮点。GS团队认为，美伊协议对欧洲利率的影响更为直接：能源缓解缩小了利率结果的分布范围，有利于降低波动率和支持主权债的carry。但他们同时强调，收益率进一步大幅下行的空间有

[... middle omitted ...]

于市场微观结构的套利机会，而非宏观判断。

*So what**: 欧洲利率市场的核心矛盾已经从“通胀是否会持续”转向了“央行会如何应对”。在这个阶段，押注方向性的大幅波动可能不如押注相对价值和曲线形态的调整。对于主权债投资者而言，欧盟债券与半核心债券之间的相对价值交易，可能比单纯做多或做空核心利率更具吸引力。

GS团队对英国利率给出了明确的偏好。他们认为，英国前端利率是G10中最有吸引力的做多标的，逻辑是通胀数据持续疲弱，加之能源价格缓解，英国央行将继续按兵不动。市场目前定价了约45个基点的加息，GS认为这一定价过于鹰派。

与此同时，英国的政治风险溢价仍然高企。Makerfield补选中Andy Burnham的大幅获胜，将继续使财政政策成为焦点。这意味着，英国国债的期限溢价难以大幅压缩。GS团队推荐2s10s英镑陡峭化交易，其逻辑是：前端利率受益于能源缓解和央行按兵不动，而长端利率则受到政治风险溢价的支撑。这是一个典型的“前端放松、长端紧绷”的曲线陡峭化策略。

\subsection{[R002] GS：市场真正低估的不是期限溢价的水平，而是它为何不再回落}
{\small\color{refgray}归类：宏观与利率：央行反应函数重塑曲线定价}

过去两年，全球债券市场一直在争论一个核心问题：长端利率的抬升，究竟是因为市场预期未来短期利率会结构性走高，还是因为投资者要求更高的“期限溢价”来补偿持有长期债券的风险？这不仅仅是学术争论。它直接决定了你的组合应该配置长债还是短债、是押注利率下行还是接受一个更高的中性利率环境。

GS在最新发布的G10期限溢价更新报告中，给出了一个比以往更锐利的分析框架。报告的核心判断不是“期限溢价现在有多高”，而是：传统基于收益率曲线推算的期限溢价模型，可能正在系统性高估当前实际期限溢价水平，因为它没能区分“结构性利率预期上移”和“周期性风险定价”。

这份报告的意义在于，它提供了一个更精确的“拆解工具”，帮助投资者判断当前债券市场定价中，哪些是可持续的结构性变化，哪些是可能均值回归的周期性溢价。对于任何管理利率风险的决策者来说，这可能是今年最重要的分析框架更新之一。

GS此次的核心贡献，是同时发布了两种期限溢价估算结果，并系统对比了它们的差异。

第一种是传统的“纯收益率曲线模型”（类似经典的ACM模型）。它只利用当前收益率曲线的形状来推断市场对未来短期利率的预期，然后将剩余部分归为期限溢价。这种模型的优点是实时性强、对市场情绪变化敏感，但缺点也很明显：它高度依赖历史均值作为长期利率预期的锚点。当经济发生结构性转变时，比如央行政策框架变化或通胀中枢上移，这种模型会倾向于把大部分收益率曲线陡峭化解释为期限溢价上升，而不是利率预期的结构性调整。

第二种是“基于调查的模型”（类似Kim-Wright方法）。它引入经济学家和交易员的短期利率预期调查数据，直接约束模型对未来利率路径的推算。这种模型的优势在于，当调查显示市场预期政策利率将结构性走高时，模型会自然地将更多长端收益率上升归因于利率预期变化，而非期限溢价。

两种模型在全球主要市场的估算结果差异显著。以美国为例，纯收益率曲线模型估算的10年期期限溢价约为110个基点，而调查模型给出的数字仅为70个基点，相差40个基点。在欧洲，这一差距约为15个基点。英国的情况更为极端：纯收益率曲线模型曾一度显示G4国家中最高的期限溢价，但调查模型却将其评为G4中最低的水平。

这意味着什么？如果调查模型更接近现实，那么当前债券市场实际上是在定价一个比传统模型所显示的要高得多的未来利率路径，而不是在定价一个异常高的风险补偿。 这对于判断长端收益率是否“超调”以及未来是否均值回归，有着根本性的不同含义。

日本市场是理解两种模型差异的最佳实验室，也是这份报告最有力的论据之一。

自2023年以来，日本调查数据中的政策利率预测持续上升，同时长期通胀预期向2\%收敛。日本国债收益率曲线显著陡峭化，伴随着日本央行资产负债表收缩和市场对久期风险的重新定价。传统收益率曲线模型会如何解读？它认为短期利率变动缓慢，基于日本长期低利率的历史锚点，模型很难在没有近期加息明确信号的情况下推高远期利率预期。因此，它把大部分收益率曲线陡峭化归因为期限溢价上升。

但调查模型给出了完全不同的答案。调查数据显示，市场参与者对日本未来2-5年的政策利率预期已经大幅上修。调查模型因此将JGB收益率上升的更大部分归因于利率预期的结构性抬升，而期限溢价的上升幅度则小得多。

GS的报告明确指出，日本是“调查模型优势的清晰例证”。如果投资者仅依赖纯收益率曲线模型，他们可能会错误地认为日本债券市场正在为“风险”定价，而实际上市场正在为“利率正常化”定价。 这两种解读对应的交易策略截然相反：前者可能建议做空长端JGB以对冲风险溢价回归，后者则可能建议做陡曲线以押注利率路径进一步上修。

这个案例对所有关注日本市场的投资者都有直接含义：日本央行退出宽松后的债券市场定价，可能并没有看起来那么“贵”，因为其中很大一部分是利率预期在追赶现实。

[... middle omitted ...]

后，调查模型显示的期限溢价上升幅度远小于纯收益率曲线模型，因为市场不仅调整了对债券供给的预期，还纳入了增长前景改善带来的实际利率结构性上移。

但纯收益率曲线模型并非没有价值。报告发现，在预测未来中长期持有久期的超额回报方面，纯收益率曲线模型表现更好。这意味着，如果投资者的目标是交易信号和短期择时，纯收益率曲线模型提供的期限溢价估计更有用；如果目标是理解市场定价的长期结构和经济含义，调查模型更可靠。

两种模型对收益率冲击的敏感性也不同。对于给定的利率水平冲击或曲线陡峭化冲击，纯收益率曲线模型的期限溢价反应幅度更大，而调查模型的反应更温和。这有助于投资者在利率波动加剧时，判断市场定价中有多少是“情绪”的过度反应，有多少是“基本面”的合理调整。

4. 一个尚未被充分回答的问题：期限溢价为何在波动率下降后仍然高企？

尽管两种模型在具体水平上存在分歧，但它们在一个关键判断上达成一致：无论用哪种方法衡量，G10国家的期限溢价在过去几年都显著上升，并处于历史区间的较高位置。

\subsection{[R003] GS：韩国市场上涨背后的结构性分歧，比指数本身更值得关注}
{\small\color{refgray}归类：区域市场：韩国结构性分歧与MSCI评估 / 风险偏好与市场情绪：权益风险指标回归中性}

KOSPI一周暴涨11\%，突破9000点，科技股领涨，外资回流。这是本周韩国市场的直观画面。

但如果只看这个数字，可能会错过这份GS周报真正想传递的信号。指数上涨的同时，资金流向、基金仓位和MSCI评估这三个维度，正在指向一个更微妙、也更需要警惕的结论：韩国市场正在经历一轮“有选择的上涨”，而非系统性重估。上涨的广度、可持续性以及背后的结构性障碍，仍然存在明显分歧。

这份报告的数据点很多，但最值得提炼的判断只有一个：韩国市场的短期动能来自事件驱动和外资回补，但中期结构性折价的修复，仍取决于MSCI评估中那些尚未解决的制度性障碍能否真正落地。 市场目前定价的是前者，而后者才是决定估值能否系统性扩张的关键变量。

本周外资净买入KOSPI市场2.5万亿韩元，推动指数大涨。单看这一数据，很容易得出“外资重新看好韩国”的结论。

但GS援引的EPFR初步数据（覆盖约46\%的AUM）给出了另一幅图景。截至2026年5月，亚洲基金和新兴市场基金对韩国股票的配置，相对于基准指数，仍然是显著低配的。亚洲基金低配195个基点，新兴市场基金低配285个基点。更关键的是，在过去一个月和过去三个月中，这两类基金都在持续削减对韩国的敞口。亚洲基金过去一个月减持105个基点，过去三个月累计减持140个基点；新兴市场基金则分别减持165和195个基点。

这意味着什么？本周的外资净买入，更可能来自交易型资金或被动型资金的短期回补，而非主动管理型基金的战略性加仓。主动基金的真实仓位仍在下降，这与指数上涨形成了方向性背离。这种背离本身就是一个需要警惕的信号：短期资金推动的上涨，如果得不到主动基金的跟进，持续性存疑。

对投资者的含义是：在跟踪韩国市场时，不能只看总量资金流向，更要区分主动与被动、交易型与配置型。当前的市场动能可能更多来自空头回补和事件驱动的投机性买盘，而非基本面逻辑驱动的系统性配置。

2. 科技硬件与半导体是外资加仓的焦点，但基金已开始减仓这一板块

本周表现最好的板块是保险（+18.4\%）、科技（+16.6\%）和零售（+5.4\%）。科技板块的盈利上调幅度也最大。这看起来是一个清晰的信号：市场在押注韩国科技股的盈利修复。

然而，GS提供的基金仓位变化数据揭示了一个更复杂的图景。在亚洲和新兴市场基金中，韩国科技硬件和半导体板块目前仍然是相对超配的（超配15个基点），但在过去三个月中，这些基金削减最多的恰恰就是科技硬件板块，减仓幅度高达60个基点。与此同时，工业板块成为过去三个月加仓最多的领域，加仓40个基点。

这组数据传达的信息是：基金正在从最拥挤、涨幅最大的科技硬件板块中撤出，转向工业等前期相对冷门的领域。这是一种典型的“轮动”而非“加仓”。科技板块本周的强劲表现，可能更多是动量交易和短期情绪推动的结果，而非机构资金的系统性增配。

对于关注韩国科技股的投资者来说，需要区分“盈利上调”和“资金流入”之间的时滞。盈利预期上调是基本面信号，但资金流向反映的是机构投资者的当前偏好。两者背离时，往往意味着市场对盈利修复的定价已经比较充分，后续上涨需要更超预期的催化剂。

这份报告最值得细读的部分，不是涨跌幅和资金流，而是MSCI市场可及性评估的详细结果。

GS明确指出，MSCI在2026年的评估中，将“投资工具可用性”从“-”上调至“+”，原因是韩国指数衍生品在国际交易所上市。这是一个积极信号，但MSCI同时保留了五个“-”的评估项：外汇自由化（缺乏完全可交割的离岸韩元）、投资者注册（IRC与LEI并存制约综合账户采用）、信息流（英文信息披露有效性待2027年全面落地）、清算与结算（仍按投资者ID结算，资金计算不透明）、以及可转移性。MSCI特别强调，其关注的是“国际机构投资者体验的实际改善”，而非宣布的改革计划。

GS的

[... middle omitted ...]

者不会。

对于资产定价的含义是：韩国市场的估值修复，不能仅仅依赖盈利增长或短期资金流入，还需要MSCI评估中那些“减号”项的真实改善。当前的市场上涨，更多是对短期事件的定价，而非对结构性折价修复的定价。

4. 报告尚未完全回答的关键问题：资金流入能否转化为持续配置

这份报告提供了丰富的数据，但有一个关键问题尚未完全回答：本周的外资净买入，到底是趋势反转的开始，还是昙花一现的脉冲？

GS的数据指向了两种可能性。乐观的一面是，韩国股市的12个月前瞻EPS本周上调了1.5\%，科技板块上调幅度最大，基本面在改善。悲观的一面是，主动基金仍在减仓，MSCI评估没有给出超预期的积极信号，韩元对美元还在贬值（本周贬值0.8\%）。

真正需要进一步验证的是：外资流入的驱动因素是什么？是美伊协议带来的风险偏好回升？是对鹰派FOMC的过度反应后的纠偏？还是韩国自身基本面改善带来的独立行情？如果是前者，那么流入的可持续性存疑；如果是后者，则需要看到更广泛的盈利上调和更持续的资金流入。

\subsection{[R004] GS：资金流向正在重塑全球资产定价的底层逻辑}
{\small\color{refgray}归类：全球资金流向：科技虹吸与新兴市场分化 / 风险偏好与市场情绪：权益风险指标回归中性}

全球资金正在经历一次结构性的再配置。这不是一次简单的板块轮动，而是一个可能影响未来12至18个月资产定价的核心变量。

GS最新一期全球资金流向周报（截至6月17日）揭示了一个清晰的信号：资金正在从分散的试探性布局，转向有明确主题的集中押注。当周全球股票基金净流入达到1260亿美元，远超此前一周的310亿美元。这一数字本身已经足够引人注目，但真正值得深入拆解的，是资金流入的结构——科技板块的吸金能力正在加速，工业板块紧随其后，而中国内地股票基金则持续面临净流出。

这份报告的核心判断是什么？不是“资金回来了”，而是“资金正在为新的宏观叙事定价”。

GS的图表清晰地展示了美国科技板块基金在最近几周的强劲净流入。这并非孤立现象。报告指出，科技基金和工业基金在全球范围内都录得了最大的净流入。但美国科技板块的流入尤为突出，其背后是AI驱动的需求增长。

这里有一个关键的分析层次。GS在报告中点明了一个经常被市场忽视的传导机制：强劲的AI需求不仅推高了增长预期，还通过推高通胀和中性利率预期，强化了美元的强势。换句话说，资金流入科技板块，表面上是在押注技术突破，实际上是在对整个宏观组合——更高的增长、更高的利率、更强的美元——进行系统性下注。

这意味着，科技板块的涨跌不再仅仅是行业基本面问题。它已经成为全球宏观资产定价的一个核心锚点。如果后续科技板块的资金流入出现逆转，其影响将远不止于纳斯达克，而是会通过美元、利率和新兴市场资金流向产生连锁反应。

在中国内地股票基金方面，GS的数据并不乐观。无论是来自中国内地本土的资金，还是来自全球其他地区的资金，都在近期出现了净流出。从累积数据看，中国内地股票基金年初至今的净流出规模在所有主要区域中排名靠前。

报告没有展开分析原因，但数据本身已经提供了足够的信息。当全球资金在科技和工业板块积极布局时，中国内地基金却持续失血，这反映出国际投资者对中国资产的风险偏好仍然偏低。值得注意的是，台湾股票基金在同一时期却录得了净流入。这种分化说明，问题不在于“中国资产”这个标签本身，而在于投资者对特定区域增长引擎和制度环境的差异化定价。

对于关注中国市场的读者而言，这里有一个需要持续跟踪的问题：中国内地基金的资金流出何时会触底？GS的数据显示，流出趋势在近几周有所加速，但尚未出现反转信号。这一趋势的逆转，可能需要更明确的宏观催化剂——无论是政策面的变化，还是企业盈利的实质性改善。

与股票市场的乐观情绪不同，固定收益市场呈现出更加谨慎的配置格局。GS报告显示，全球固定收益基金继续获得资金支持，但资金偏好明显偏向短期。短久期债券基金和通胀保护债券基金持续获得净流入。

这是一个经典的“防御性配置”信号。当投资者同时买入股票（尤其是科技股）和短久期债券时，他们实际上是在表达一个矛盾的预期：一方面看好AI驱动的增长前景，另一方面又对利率路径和通胀前景感到不确定。这种“股债双买”的格局，在历史上往往出现在宏观叙事切换的过渡期。

报告还提到，新兴市场硬通货债券基金获得净流入，而本币债券基金则遭遇净流出。这进一步印证了投资者对新兴市场货币的谨慎态度——他们愿意承担信用风险，但不愿意承担汇率风险。这和GS在报告中提到的美元净需求强劲、人民币净流出最大的趋势完全一致。

GS这份周报提供了丰富的流量数据，但有一个关键问题没有完全回答：这些资金流入的可持续性如何？

从历史经验看，科技板块的集中流入往往伴随着较高的波动风险。一旦AI相关企业的盈利预期出现修正，或者宏观数据出现意外波动，这些资金可能会迅速撤离。GS的图表显示，美国科技板块的流入在最近几周确实很强劲，但类似的集中流入在历史上并非没有出现过，而后续往往伴随着资金流向的剧烈反转。

另一个未解问题是：欧洲资金的流入是否可持续？报告中提到，欧元区股

[... middle omitted ...]

工业双双走强，这可能意味着AI的投资逻辑正在从“软件和芯片”向“硬件和应用”扩散。这种扩散一旦确认，将改变整个周期性板块的定价框架。

第二，美元强势与新兴市场资金流向的负反馈循环。GS的报告明确指出，AI驱动的需求推高了美国的中性利率，进而强化了美元。而美元的强势又导致新兴市场（尤其是中国）的资本外流。这个循环何时会自我修正？只有当美国经济数据出现明显放缓，或者新兴市场的增长预期出现实质性改善时，这个循环才有可能被打破。

第三，固定收益市场的短久期偏好何时转向。如果投资者开始从短久期债券转向长久期债券，那将是一个重要的信号——意味着市场对利率路径的信心开始恢复。GS的数据目前还没有显示这种转向，但这是一个需要持续监控的指标。

GS的周报提供了丰富的流量数据，但受到周报格式的限制，很多重要的二阶效应并没有被充分展开。例如，资金流入科技板块对中小盘科技股的溢出效应如何？中国内地基金流出的同时，中国债券市场的资金流向又是怎样的？跨境资金流动对汇率定价的短期冲击有多大？

\subsection{[R005] GS：市场低估的不是资本开支本身，而是资本开支正在重塑行业竞争格局}
{\small\color{refgray}归类：AI与算力：资本开支重塑行业竞争格局 / 能源与大宗：资本开支回暖与结构性变化}

GS：市场低估的不是资本开支本身，而是资本开支正在重塑行业竞争格局

这份GS研报传递了一个核心信号：全球正在进入一个由资本开支驱动的“后现代周期”，但市场对它的理解仍然停留在“谁在花钱”的层面，而忽略了“钱花在哪里、谁有能力花、花完之后谁受益”这三个更深层的结构性变化。

研报覆盖了从AI数据中心、油气勘探到电力基础设施的广泛领域，但将这些分散的线索串联起来，你会发现一个共同的主线：资本开支的回报逻辑正在从“规模扩张”转向“议价权重构”。那些能够将资本投入转化为行业准入壁垒、技术标准定义权或稀缺资源控制权的企业，才是这个周期真正的赢家。

以下是我们从这份报告中提炼出的五个关键洞察，以及一个尚未被充分讨论的隐忧。

1. 资本开支周期的本质不是“花钱”，而是“重新划定谁有资格进场”

GS全球策略师Peter Oppenheimer在研报中提出了一个关键判断：当前资本开支周期与以往不同，其回报更多来自每股盈利增长而非估值扩张，这意味着个股之间的分化将显著加大，主动选股的空间正在打开。

这个判断的潜台词是：市场正在从“水涨船高”的阶段进入“优胜劣汰”的阶段。过去几年，低利率环境下，几乎所有企业都能轻松融资、扩大投资，资本开支的回报差异被流动性掩盖。而现在，随着利率维持高位、融资成本上升，资本开支的“准入门槛”正在提高——只有那些拥有稳定现金流、明确战略方向和技术壁垒的企业，才能持续投入并看到回报。

以数据中心为例。GS欧洲公用事业分析师Alberto Gandolfi的最新估算显示，欧洲数据中心管道已接近500GW，相当于当前欧盟电力需求的1.5倍，六个月前这个数字还只有约300GW。这意味着，谁能在未来几年获得稳定的电力供应、谁能在电网接入和土地审批上占据先机，谁就掌握了AI基础设施的“入场券”。这不是简单的“花钱建数据中心”的问题，而是“谁有资格建、建在哪里、能否持续运营”的竞争。

同样的逻辑也适用于油气勘探。GS能源团队注意到，即使布伦特油价回落至80美元/桶附近，国际石油公司和国有石油公司仍在加大前沿勘探投入。原因很简单：过去十年，全球油气储采比下降了60\%，而美国页岩油这个曾经压制油价的“供应缓冲器”正在成熟。在能源安全和储量焦虑的双重压力下，勘探投入不再是可选项，而是必选项。但能够承担这种长期、高风险资本开支的，只有那些资产负债表稳健、技术能力领先的头部企业。

所以，资本开支周期的核心不是“花多少钱”，而是“谁在花，以及花完之后他们能获得什么样的竞争位置”。

2. AI资本开支的“暗工厂”效应：软件行业的成本结构正在被重写

研报中一个常被忽视的细节来自GS美国科技团队组织的一场专家电话会。专家描述了一个“暗工厂”模型：编码代理自动编写代码并部署到生产环境，站点可靠性工程代理评估生产行为、诊断问题，并将信息反馈回开发循环，从而在软件创建和软件运营之间建立更紧密的闭环。

这个描述的意义远远超出了“AI提高效率”的泛泛之谈。它指向的是一个更根本的变化：软件行业的成本结构正在从“人力密集型”转向“资本密集型”。过去，软件公司的核心资产是工程师的脑力和时间，资本开支主要集中在服务器和办公空间。而“暗工厂”模型意味着，软件开发和运维的绝大部分环节可以被自动化，企业的竞争优势将从“谁能雇到更多更好的工程师”转向“谁能构建出更高效、更可靠的自动化系统”。

这对软件公司的商业模式意味着什么？GSIT支出调查显示，目前88\%的CIO预计AI支出占IT预算的比例不到10\%，但三年后，42\%的CIO预计生成式AI支出将超过10\%。这组数据说明，企业仍在测试和验证AI的实际价值，但一旦验证通过，资本开支将加速涌入。那些能够提供“AI原生基础设施”的企业——无论是芯片、服务器、数据中心还是自动化运维平台——将获

[... middle omitted ...]

S分析师Michele Della Vigna与Viridien（前CGG）CFO的交流中得到的反馈是：“前沿勘探回来了。”

这不仅仅是地缘政治事件的短期反应。研报明确指出，油气巨头对勘探投入的关注在近期中东冲突之前就已经开始回升。原因在于，投资者开始奖励勘探支出而非仅仅回购股票，因为储采比在过去十年下降了60\%，而美国页岩油这个曾经压制油价的供应来源正在成熟。能源安全只是加剧了这一趋势。

GS预计，这些勘探项目在60美元/桶的油价环境下就能实现盈利，而当前布伦特油价仍在80美元/桶附近。这意味着，即使油价出现较大幅度的回落，这些项目的经济性仍然成立。因此，看空油价需要非常极端的假设才能推翻这一投资逻辑。

但市场对此的定价可能仍不充分。油气服务板块在油价回落后表现疲弱，投资者似乎仍在用“油价决定一切”的旧框架来评估这些公司。然而，如果勘探支出的结构性增长成立，那么油气服务公司的收入来源将从“跟随油价波动”转向“受益于资本开支的长期上行”，其估值逻辑需要重新审视。

\section{Disclaimer}
{\scriptsize\color{refgray}
This community is a paid learning and information-sharing community created for general educational purposes only. The membership fee is charged solely for access to community discussions, educational content, organization, and operational costs. It is not an advisory fee, management fee, performance fee, brokerage fee, or compensation for personalized financial, investment, legal, tax, or professional advice.\par
I participate and share information solely as an individual and for educational discussion among community members. I am not acting as a financial adviser, investment adviser, broker, dealer, portfolio manager, legal adviser, tax adviser, accountant, fiduciary, or any other licensed professional adviser. Nothing shared in this community constitutes, and should not be interpreted as, financial advice, investment advice, legal advice, tax advice, a recommendation, solicitation, offer, endorsement, instruction, or guarantee to buy, sell, hold, short, trade, or invest in any security, cryptocurrency, fund, derivative, strategy, or other financial product.\par
All content shared in this community, including messages, discussions, links, files, charts, examples, market commentary, watchlists, AI-generated or AI-polished summaries, and third-party materials, is general, impersonal, and educational in nature. It does not take into account any individual's financial situation, investment objectives, risk tolerance, investment horizon, liquidity needs, tax status, legal restrictions, or suitability requirements. No content should be relied upon as suitable for any specific person, account, portfolio, or financial decision.\par
Information shared in this community may be incomplete, inaccurate, outdated, speculative, AI-generated, AI-edited, or based on third-party internet sources that have not been independently verified. I make no representation or warranty as to the accuracy, completeness, timeliness, reliability, or suitability of any information shared.\par
Financial markets involve substantial risk, including the possible loss of principal. Past performance, historical data, hypothetical examples, simulated results, backtests, personal opinions, or educational examples do not guarantee future results. Each member is solely responsible for conducting their own research, exercising independent judgment, and making their own decisions. Before making any financial, investment, legal, or tax decision, members should consult qualified and licensed professionals.\par
Participation in this community, payment of any membership fee, or communication with me or other members does not create any adviser-client, fiduciary, professional, contractual, agency, partnership, employment, or other advisory relationship. By joining or participating in this community, you acknowledge and agree that you are solely responsible for your own decisions, actions, transactions, profits, losses, risks, and consequences, and that I shall not be liable for any direct, indirect, incidental, consequential, financial, legal, tax, or other loss or damage arising from or related to any content, discussion, information, or material shared in this community.\par
}
\end{document}
